% ============================================================
%  Chapter 3 – Livestock enteric methane emissions
% ============================================================
\chapter{Livestock enteric methane emissions}
\label{ch:enteric}

\textit{Authors: A.W.\ Alemu and S.J.\ Pogue}

% ─────────────────────────────────────────────────────────────────────────────
\section{Beef cattle}
\label{sec:beef}

Management data are input for each management period. Emissions are calculated
on a daily basis.

\noindent\textbf{Assumptions:}
\begin{itemize}[leftmargin=1.5em]
    \item Cattle feed intake is equal to energy requirements
    \item All feed is utilized --- waste feed emissions are not accounted for
    \item All cows are pregnant and 5\,kg of protein is retained for every pregnancy
    \item Calves consume 1\% of their own body weight as solid food
    \item Calves retain 20\% of protein intake from dry feed and 40\% of protein
          intake from milk
    \item Enteric CH$_4$ emissions from calves are calculated using the IPCC $Y_m$
          approach only (Section~\ref{subsec:beef-calves})
\end{itemize}

% ·····················································
\subsection{IPCC \texorpdfstring{$Y_m$}{Ym} approach for beef cattle (except calves)}
\label{subsec:beef-ym}

Enteric CH$_4$ calculations should be completed for each cattle group (except calves).

\begin{equation}
    \mathrm{Weight} = \mathrm{ADG} \times \mathrm{days} + \mathrm{wt}_{\mathrm{initial}}
    \tag{Eq.\ 3.1.1.1}
\end{equation}

\noindent where
\begin{description}[font=\normalfont, leftmargin=9em, style=sameline, nosep]
    \item[$\mathrm{Weight}$] Weight (kg~head$^{-1}$)
    \item[$\mathrm{ADG}$] Average daily gain (kg~head$^{-1}$~day$^{-1}$), calculated
          using Eq.~3.1.1.7
    \item[$\mathrm{days}$] Total number of days since start of management period
    \item[$\mathrm{wt}_{\mathrm{initial}}$] Initial weight (kg~head$^{-1}$) (Table~16,
          by beef cattle group, or user-defined)
\end{description}

\bigskip
\noindent\textbf{Net energy requirements}

$C_{f,\mathrm{adjusted}}$ represents the influence of temperature on the maintenance
coefficient ($C_f$) and varies by day due to daily average temperature.

\begin{equation}
    C_{f,\mathrm{adjusted}} = C_f + \bigl(0.0048 \times (20 - \mathrm{Temperature})\bigr)
    \tag{Eq.\ 3.1.1.2}
\end{equation}

\noindent IPCC (2019), Eq.~10.2

\noindent where
\begin{description}[font=\normalfont, leftmargin=9em, style=sameline, nosep]
    \item[$C_{f,\mathrm{adjusted}}$] Maintenance coefficient adjusted for temperature
          (MJ~day$^{-1}$~kg$^{-1}$)
    \item[$C_f$] Baseline maintenance coefficient (MJ~day$^{-1}$~kg$^{-1}$) (Table~16,
          by beef cattle group)
    \item[$\mathrm{Temperature}$] Average daily temperature (upper limit = 20\textdegree{}C;
          if temperature $>$ 20\textdegree{}C or if cattle are housed in a barn, use
          20\textdegree{}C)
\end{description}

\noindent\textit{Note: While this is calculated and labelled as $C_f$ in the model
code, it is labelled ``Maintenance coefficient'' in the model interface. By default,
daily Temperature data are obtained from the NASA POWER data access viewer
(\url{https://power.larc.nasa.gov/data-access-viewer/}). The model user can override
these data using custom daily climate data. In the event that neither NASA daily climate
data nor user-defined daily climate data are available, Holos uses 1990--2017 monthly
climate normals for each SLC polygon as a default.}

\begin{equation}
    \mathrm{NE}_{\mathrm{maintenance}} =
        C_{f,\mathrm{adjusted}} \times \mathrm{Weight}^{0.75}
    \tag{Eq.\ 3.1.1.3}
\end{equation}

\noindent IPCC (2019), Eq.~10.3

\noindent where
\begin{description}[font=\normalfont, leftmargin=9em, style=sameline, nosep]
    \item[$\mathrm{NE}_{\mathrm{maintenance}}$] Net energy for maintenance
          (MJ~head$^{-1}$~day$^{-1}$)
\end{description}

\begin{equation}
    \mathrm{NE}_{\mathrm{activity}} = C_a \times \mathrm{NE}_{\mathrm{maintenance}}
    \tag{Eq.\ 3.1.1.4}
\end{equation}

\noindent IPCC (2019), Eq.~10.4

\noindent where
\begin{description}[font=\normalfont, leftmargin=9em, style=sameline, nosep]
    \item[$\mathrm{NE}_{\mathrm{activity}}$] Net energy for activity
          (MJ~head$^{-1}$~day$^{-1}$)
    \item[$C_a$] Feeding activity coefficient (Table~17, by activity type)
\end{description}

\noindent For lactating beef cows only (use only when cows are lactating ---
when \#calves $>$ 0):

\begin{equation}
    \mathrm{NE}_{\mathrm{lactation}} =
        \bigl[\mathrm{milk}_{\mathrm{prod}} \times 1.47
            + 0.40 \times \mathrm{fat}_{\mathrm{content}}\bigr]
        \times \frac{\#\mathrm{calves}}{\#\mathrm{cows}}
    \tag{Eq.\ 3.1.1.5}
\end{equation}

\noindent IPCC (2019), Eq.~10.8

\noindent where
\begin{description}[font=\normalfont, leftmargin=9em, style=sameline, nosep]
    \item[$\mathrm{NE}_{\mathrm{lactation}}$] Net energy for lactation
          (MJ~head$^{-1}$~day$^{-1}$)
    \item[$\mathrm{milk}_{\mathrm{prod}}$] Milk production (kg~head$^{-1}$~day$^{-1}$)
    \item[$\mathrm{fat}_{\mathrm{content}}$] Fat content (\%)
    \item[$\#\mathrm{calves}$] Number of calves
    \item[$\#\mathrm{cows}$] Number of cows
\end{description}

\noindent Holos V4 uses the following constant values:
$\mathrm{milk}_{\mathrm{prod}} = 8$\,(kg~head$^{-1}$~day$^{-1}$),
$\mathrm{fat}_{\mathrm{content}} = 4$\,(\%).

\noindent\textit{Note: \#calves/\#cows adjusts for cows without calves and
averages energy for lactation across all individuals.}

\noindent For pregnant beef cows only:

\begin{equation}
    \mathrm{NE}_{\mathrm{pregnancy}} = 0.10 \times \mathrm{NE}_{\mathrm{maintenance}}
    \tag{Eq.\ 3.1.1.6}
\end{equation}

\noindent IPCC (2019), Eq.~10.13

\noindent where
\begin{description}[font=\normalfont, leftmargin=9em, style=sameline, nosep]
    \item[$\mathrm{NE}_{\mathrm{pregnancy}}$] Net energy for pregnancy
          (MJ~head$^{-1}$~day$^{-1}$)
    \item[$0.1$] Pregnancy coefficient --- IPCC (2019) default for cattle and buffalo
\end{description}

\noindent\textit{This equation averages pregnancy energy requirements over the
entire year.}

\bigskip
\noindent\textbf{Average daily gain (ADG), net energy for gain}

For all beef cattle groups, default ADG values are provided for each management period,
based on default start and end weights and the number of days in each management period.
As for most other inputs in Holos, this information can be overridden with user-defined
inputs.

\begin{equation}
    \mathrm{ADG} = \frac{\mathrm{wt}_{\mathrm{final}} - \mathrm{wt}_{\mathrm{initial}}}
                        {\#\mathrm{days}}
    \tag{Eq.\ 3.1.1.7}
\end{equation}

\noindent where
\begin{description}[font=\normalfont, leftmargin=9em, style=sameline, nosep]
    \item[$\mathrm{ADG}$] Average daily gain (kg~head$^{-1}$~day$^{-1}$)
    \item[$\mathrm{wt}_{\mathrm{final}}$] Final weight (kg~head$^{-1}$) (Table~16,
          by beef cattle group)
    \item[$\mathrm{wt}_{\mathrm{initial}}$] Initial weight (kg~head$^{-1}$) (Table~16,
          by beef cattle group)
    \item[$\#\mathrm{days}$] Number of days --- includes all months where animals
          $>0$ are entered
\end{description}

\noindent For all beef animals (excluding calves):

\begin{equation}
    \mathrm{NE}_{\mathrm{gain}} = 22.02 \times
        \frac{\mathrm{Weight}}{C_d \times \mathrm{wt}_{\mathrm{final}}^{0.75}}
        \times \mathrm{ADG}^{1.097}
    \tag{Eq.\ 3.1.1.8}
\end{equation}

\noindent IPCC (2019), Eq.~10.6

\noindent where
\begin{description}[font=\normalfont, leftmargin=9em, style=sameline, nosep]
    \item[$\mathrm{NE}_{\mathrm{gain}}$] Net energy for gain (MJ~head$^{-1}$~day$^{-1}$)
    \item[$\mathrm{Weight}$] Weight (kg~head$^{-1}$)
    \item[$\mathrm{wt}_{\mathrm{final}}$] Final weight of animal (kg) (Table~16)
    \item[$C_d$] Gain coefficient (Table~16, by cattle group)
    \item[$\mathrm{ADG}$] Average daily gain (kg~head$^{-1}$~day$^{-1}$)
\end{description}

\bigskip
\noindent\textbf{Ratios of net energy available to digestible energy}

\begin{equation}
    \mathrm{REM} = 1.123 - 4.092 \times 10^{-3} \times \mathrm{TDN}
        + 1.126 \times 10^{-5} \times \mathrm{TDN}^{2}
        - \frac{25.4}{\mathrm{TDN}}
    \tag{Eq.\ 3.1.1.9}
\end{equation}

\noindent IPCC (2019), Eq.~10.14

\noindent where
\begin{description}[font=\normalfont, leftmargin=6em, style=sameline, nosep]
    \item[$\mathrm{REM}$] Ratio of net energy available in diet for maintenance to
          digestible energy consumed
    \item[$\mathrm{TDN}$] Percent total digestible nutrients in feed (Table~18,
          by diet)
\end{description}

\noindent\textit{Note: TDN is to be entered as a percentage (e.g.\ as 81 not 0.81).}

\begin{equation}
    \mathrm{REG} = 1.164 - 5.160 \times 10^{-3} \times \mathrm{TDN}
        + 1.308 \times 10^{-5} \times \mathrm{TDN}^{2}
        - \frac{37.4}{\mathrm{TDN}}
    \tag{Eq.\ 3.1.1.10}
\end{equation}

\noindent IPCC (2019), Eq.~10.15

\noindent where
\begin{description}[font=\normalfont, leftmargin=6em, style=sameline, nosep]
    \item[$\mathrm{REG}$] Ratio of net energy available in diet for gain to
          digestible energy consumed
\end{description}

\noindent\textit{Note: TDN is to be entered as a percentage (e.g.\ as 81 not 0.81).}

\bigskip
\noindent\textbf{Gross energy intake}

\begin{multline}
    \mathrm{GEI} =
        \left(\frac{\mathrm{NE}_{\mathrm{maintenance}} + \mathrm{NE}_{\mathrm{activity}}
            + \mathrm{NE}_{\mathrm{lactation}} + \mathrm{NE}_{\mathrm{pregnancy}}}
            {\mathrm{REM}}
        + \frac{\mathrm{NE}_{\mathrm{gain}}}{\mathrm{REG}}\right)
        \times \frac{100}{\mathrm{TDN}}
    \tag{Eq.\ 3.1.1.11}
\end{multline}

\noindent IPCC (2019), Eq.~10.16

\noindent where
\begin{description}[font=\normalfont, leftmargin=6em, style=sameline, nosep]
    \item[$\mathrm{GEI}$] Gross energy intake (MJ~head$^{-1}$~day$^{-1}$)
\end{description}

\noindent\textit{Note: TDN is to be entered as a percentage (e.g.\ as 81 not 0.81).}

\bigskip
\noindent\textbf{Enteric CH$_4$ emission}

\begin{equation}
    \mathrm{CH}_{4,\mathrm{enteric\text{-}rate}} =
        \mathrm{GEI} \times \frac{Y_m}{55.65}
        \times \left(1 - \frac{\mathrm{AR}}{100}\right)
    \tag{Eq.\ 3.1.1.12}
\end{equation}

\noindent Derived from IPCC (2019), Eq.~10.21

\noindent where
\begin{description}[font=\normalfont, leftmargin=9em, style=sameline, nosep]
    \item[$\mathrm{CH}_{4,\mathrm{enteric\text{-}rate}}$] Enteric CH$_4$ emission rate
          (kg~head$^{-1}$~day$^{-1}$)
    \item[$Y_m$] Methane conversion factor (Table~18, by diet)
    \item[$55.65$] Energy content of CH$_4$ (MJ~kg$^{-1}$~CH$_4$)
    \item[$\mathrm{AR}$] Additive reduction factor (Table~19, by additive)
\end{description}

\noindent\textit{Note: When a custom diet is created by the model user, $Y_m$ is first
determined on the basis of the forage content --- if forage content $<$ 15\%, a default
$Y_m$ value of 0.04 is used. If forage content $>$ 15\% then $Y_m$ is calculated based
on the TDN content of the diet.}

\begin{equation}
    \mathrm{CH}_{4,\mathrm{enteric}} =
        \mathrm{CH}_{4,\mathrm{enteric\text{-}rate}} \times \#\mathrm{cattle}
    \tag{Eq.\ 3.1.1.13}
\end{equation}

\noindent where
\begin{description}[font=\normalfont, leftmargin=9em, style=sameline, nosep]
    \item[$\mathrm{CH}_{4,\mathrm{enteric}}$] Enteric CH$_4$ emissions by animal group
          (kg~CH$_4$~day$^{-1}$)
    \item[$\#\mathrm{cattle}$] Number of beef cattle
\end{description}

\bigskip
\noindent\textbf{Other enteric CH$_4$ estimates}

In addition to using the default IPCC (2019) $Y_m$ conversion factors to estimate
enteric CH$_4$ emissions from beef cattle, the inclusion of different mathematical
equations (linear and exponential) with varying input requirements provides the model
user (and in particular researchers) with additional options. For beef cattle, the
selected equations were developed and evaluated by van Lingen et al.\ (2019). Four
equations were selected and divided into two groups based on the concentration of forage
in the diet. Two equations are for diets containing $\geq 25$\% forage DM (high forage
diets, Eq.~3.1.1.14 and Eq.~3.1.1.15) and the remaining two equations are for diets
containing $< 25$\% forage DM (low forage diets, Eq.~3.1.1.16 and Eq.~3.1.1.17).

\noindent\textit{High forage diet alternatives ($\geq 25$\% dietary forage)}

\begin{multline}
    \mathrm{CH}_{4,\mathrm{enteric}(\mathrm{B1})\mathrm{\text{-}rate}} =
        \left[-35.0 + 0.08 \times \mathrm{BW} + 1.2 \times \mathrm{Forage}
        - 69.8 \times \mathrm{EEI}^{3}
        + 3.14 \times \frac{\mathrm{GEI}}{4.184}\right]
        \times \frac{1}{1000}
    \tag{Eq.\ 3.1.1.14}
\end{multline}

\noindent Escobar-Bahamondes et al.\ (2017), Eq.~AL-OR

\noindent where
\begin{description}[font=\normalfont, leftmargin=9em, style=sameline, nosep]
    \item[$\mathrm{CH}_{4,\mathrm{enteric}(\mathrm{B1})\mathrm{\text{-}rate}}$]
          Enteric CH$_4$ emissions (kg~head$^{-1}$~day$^{-1}$)
    \item[$\mathrm{BW}$] Body weight (kg)
    \item[$\mathrm{Forage}$] \% forage in the diet (\% of DM)
    \item[$\mathrm{EEI}$] Ether extract/crude fat intake (kg~head$^{-1}$~day$^{-1}$)
    \item[$\mathrm{GEI}$] Gross energy intake (MJ~head$^{-1}$~day$^{-1}$)
\end{description}

\noindent Holos V4 uses the following calculation:
$\mathrm{EEI} = \mathrm{DMI}\,(\text{kg/day}) \times [\mathrm{EE\ concentration\ in\ feed\ or\ diet}\ (\%\mathrm{DM})/100]$
$+ \mathrm{DMI}\,(\text{kg~day}^{-1}) \times (\%\text{Fat supplement added as diet additives}/100)$

\begin{equation}
    \mathrm{CH}_{4,\mathrm{enteric}(\mathrm{B2})\mathrm{\text{-}rate}} =
        \left(-6.41 + 11.3 \times \mathrm{DMI}
            + 0.557 \times \mathrm{Forage}
            + 0.0996 \times \mathrm{BW}\right)
        \times \frac{1}{1000}
    \tag{Eq.\ 3.1.1.15}
\end{equation}

\noindent van Lingen et al.\ (2019), Eq.~17

\noindent where
\begin{description}[font=\normalfont, leftmargin=9em, style=sameline, nosep]
    \item[$\mathrm{CH}_{4,\mathrm{enteric}(\mathrm{B2})\mathrm{\text{-}rate}}$]
          Enteric CH$_4$ emission rate (kg~head$^{-1}$~day$^{-1}$)
    \item[$\mathrm{DMI}$] DM intake (kg~head$^{-1}$~day$^{-1}$), calculated using
          Eq.~11.3.1.1
    \item[$\mathrm{Forage}$] Dietary forage (\% of DM)
    \item[$\mathrm{BW}$] Body weight (kg)
\end{description}

\noindent\textit{Low forage diet alternatives ($< 25$\% dietary forage)}

\begin{multline}
    \mathrm{CH}_{4,\mathrm{enteric}(\mathrm{B3})\mathrm{\text{-}rate}} =
        \left[-10.1 + 0.21 \times \mathrm{BW}
            + 0.36 \times \mathrm{DMI}^{2}
            - 69.2 \times \mathrm{EEI}^{3}
            + 13 \times \frac{\mathrm{CPI}}{\mathrm{NDFI}}
            - 4.9 \times \frac{\mathrm{StarchI}}{\mathrm{NDFI}}\right]
        \times \frac{1}{1000}
    \tag{Eq.\ 3.1.1.16}
\end{multline}

\noindent Escobar-Bahamondes et al.\ (2017), Eq.~LF-MC

\noindent where
\begin{description}[font=\normalfont, leftmargin=7em, style=sameline, nosep]
    \item[$\mathrm{CH}_{4,\mathrm{enteric}(\mathrm{B3})\mathrm{\text{-}rate}}$]
          Enteric CH$_4$ emission rate (kg~head$^{-1}$~day$^{-1}$)
    \item[$\mathrm{BW}$] Body weight (kg)
    \item[$\mathrm{DMI}$] DM intake (kg~head$^{-1}$~day$^{-1}$), calculated using
          Eq.~11.3.1.1
    \item[$\mathrm{EEI}$] Ether extract/crude fat intake (kg~head$^{-1}$~day$^{-1}$)
    \item[$\mathrm{CPI}$] Crude protein intake (kg~head$^{-1}$~day$^{-1}$)
    \item[$\mathrm{NDFI}$] Neutral detergent fiber intake (kg~head$^{-1}$~day$^{-1}$)
    \item[$\mathrm{StarchI}$] Starch intake (kg~head$^{-1}$~day$^{-1}$)
\end{description}

\noindent Holos V4 uses the following calculations:
$\mathrm{CPI} = \mathrm{DMI}\,(\text{kg/day}) \times [\mathrm{CP\ concentration\ in\ feed\ or\ diet}\ (\%\mathrm{DM})/100]$;
$\mathrm{NDFI} = \mathrm{DMI}\,(\text{kg/day}) \times [\mathrm{NDF\ concentration\ in\ feed\ or\ diet}\ (\%\mathrm{DM})/100]$;
$\mathrm{StarchI} = \mathrm{DMI}\,(\text{kg/day}) \times [\mathrm{Starch\ concentration\ in\ feed\ or\ diet}\ (\%\mathrm{DM})/100]$

\begin{equation}
    \mathrm{CH}_{4,\mathrm{enteric}(\mathrm{B4})\mathrm{\text{-}rate}} =
        \left[48.2 + 14.4 \times \mathrm{DMI}
            - 20.5 \times \frac{\mathrm{Starch}}{\mathrm{NDF}}\right]
        \times \frac{1}{1000}
    \tag{Eq.\ 3.1.1.17}
\end{equation}

\noindent Ellis et al.\ (2009), Eq.~N

\noindent where
\begin{description}[font=\normalfont, leftmargin=9em, style=sameline, nosep]
    \item[$\mathrm{CH}_{4,\mathrm{enteric}(\mathrm{B4})\mathrm{\text{-}rate}}$]
          Enteric CH$_4$ emissions rate (kg~head$^{-1}$~day$^{-1}$)
    \item[$\mathrm{DMI}$] DM intake (kg~head$^{-1}$~day$^{-1}$), calculated using
          Eq.~11.3.1.1
    \item[$\mathrm{Starch}$] Dietary starch (\% of DM)
    \item[$\mathrm{NDF}$] Dietary neutral detergent fiber (\% of DM)
\end{description}

\begin{equation}
    \mathrm{CH}_{4,\mathrm{enteric}(n)} =
        \mathrm{CH}_{4,\mathrm{enteric}(n)\mathrm{\text{-}rate}}
        \times \#\mathrm{cattle}
    \tag{Eq.\ 3.1.1.18}
\end{equation}

\noindent where
\begin{description}[font=\normalfont, leftmargin=9em, style=sameline, nosep]
    \item[$\mathrm{CH}_{4,\mathrm{enteric}(n)}$] Enteric CH$_4$ emissions by animal
          group (kg~CH$_4$~day$^{-1}$), $n = \mathrm{B1, B2, B3, B4}$
    \item[$\#\mathrm{cattle}$] Number of cattle
\end{description}

% ·····················································
\subsection{Enteric CH\texorpdfstring{$_4$}{4} emissions from beef calves}
\label{subsec:beef-calves}

Enteric CH$_4$ emissions from beef calves are estimated based on a simplified Tier~2
approach from IPCC (2019) for the estimation of daily DMI. Enteric CH$_4$ emissions for
this animal group are estimated only for those days when the calves are on a non-milk
(i.e., forage) diet. To calculate enteric CH$_4$ from these calves, we assume that, in
addition to consuming milk, calves also graze on forage from 01 July to 31 October
(123 days, according to Legesse et al.\ 2016, 2018) on the same pasture (native or tame)
as the cows. Prior to that, the calves consume milk only and thus generate no enteric
CH$_4$.

\begin{equation}
    \mathrm{Weight} = \mathrm{ADG} \times \mathrm{days} + \mathrm{wt}_{\mathrm{initial}}
    \tag{Eq.\ 3.1.2.1}
\end{equation}

\noindent where
\begin{description}[font=\normalfont, leftmargin=9em, style=sameline, nosep]
    \item[$\mathrm{Weight}$] Weight (kg~head$^{-1}$)
    \item[$\mathrm{ADG}$] Average daily gain (kg~head$^{-1}$~day$^{-1}$), calculated
          according to Eq.~3.1.1.7
    \item[$\mathrm{days}$] Total number of days since start of management period
    \item[$\mathrm{wt}_{\mathrm{initial}}$] Initial weight (kg~head$^{-1}$) (Table~16,
          by beef cattle group)
\end{description}

\begin{equation}
    \mathrm{DMI} = \frac{\mathrm{Weight}^{0.75}
        \times \bigl(0.0582 \times \mathrm{NE}_{\mathrm{mf}}
            - 0.00266 \times \mathrm{NE}_{\mathrm{mf}}^{2}
            - 0.1128\bigr)}
        {0.239 \times \mathrm{NE}_{\mathrm{mf}}}
    \tag{Eq.\ 3.1.2.2}
\end{equation}

\noindent Derived from IPCC (2019), Eq.~10.17

\noindent where
\begin{description}[font=\normalfont, leftmargin=7em, style=sameline, nosep]
    \item[$\mathrm{DMI}$] DM intake (kg~head$^{-1}$~day$^{-1}$)
    \item[$\mathrm{Weight}$] Weight (kg~head$^{-1}$)
    \item[$\mathrm{NE}_{\mathrm{mf}}$] Estimated dietary net energy concentration of
          diet or default values (MJ~kg$^{-1}$) (Table~20)
\end{description}

\noindent\textit{Note: For beef and dairy calves on a milk-only diet, the DMI is
estimated as: $\mathrm{DMI} = \mathrm{Weight} \times 0.01$}

\begin{equation}
    \mathrm{GEI} = \mathrm{DMI} \times 18.45
    \tag{Eq.\ 3.1.2.3}
\end{equation}

\noindent IPCC (2019)

\noindent where
\begin{description}[font=\normalfont, leftmargin=6em, style=sameline, nosep]
    \item[$\mathrm{GEI}$] Gross energy intake (MJ~head$^{-1}$~day$^{-1}$)
    \item[$18.45$] Conversion factor for gross energy per kg of DM (MJ~kg$^{-1}$)
\end{description}

\noindent Use Eq.~3.1.1.12 and Eq.~3.1.1.13 to calculate enteric CH$_4$ emissions
for beef calves.

% ·····················································
\subsection{Total enteric CH\texorpdfstring{$_4$}{4} emissions}
\label{subsec:beef-total}

Enteric CH$_4$ emissions should be summed for all relevant management periods and
beef cattle groups.

\begin{equation}
    \mathrm{Total\_CH}_{4,\mathrm{enteric}} =
        \sum_{\mathrm{all\ scenario\ cattle}}
        \mathrm{CH}_{4,\mathrm{enteric}}
    \tag{Eq.\ 3.1.3.1}
\end{equation}

\noindent IPCC (2019), Eq.~10.20

\noindent where
\begin{description}[font=\normalfont, leftmargin=9em, style=sameline, nosep]
    \item[$\mathrm{Total\_CH}_{4,\mathrm{enteric}}$] Total enteric CH$_4$ emissions
          from beef cattle (kg~CH$_4$~year$^{-1}$)
    \item[$\mathrm{CH}_{4,\mathrm{enteric}}$] Enteric CH$_4$ emissions (kg~CH$_4$)
\end{description}

% ─────────────────────────────────────────────────────────────────────────────
\section{Dairy cattle}
\label{sec:dairy}

Management data are input for each management period. Emissions are calculated
on a daily basis.

\noindent\textbf{Assumptions:}
\begin{itemize}[leftmargin=1.5em]
    \item All cows are pregnant
    \item 5\,kg of protein is retained for every pregnancy
    \item All feed is utilized --- waste feed emissions are not accounted for
    \item Cattle feed intake is equal to energy requirements
    \item Three lactation periods each year are assumed, and milk production can be
          defined separately for each of these periods
    \item A single default management period for dairy calves of 3 months is assumed,
          during which time these animals are milk-fed only (after this, emissions for
          weaned calves may be calculated under a `Dairy heifers' group)
    \item For all milk-fed-only dairy calves, enteric CH$_4$ emissions are assumed to
          be zero
    \item Manure must be applied at least once per year
\end{itemize}

% ·····················································
\subsection{IPCC \texorpdfstring{$Y_m$}{Ym} approach for dairy cattle (except calves)}
\label{subsec:dairy-ym}

Enteric CH$_4$ calculations should be completed for each dairy cattle group (except calves).

\begin{equation}
    \mathrm{Weight} = \mathrm{ADG} \times \mathrm{days} + \mathrm{wt}_{\mathrm{initial}}
    \tag{Eq.\ 3.2.1.1}
\end{equation}

\noindent where
\begin{description}[font=\normalfont, leftmargin=9em, style=sameline, nosep]
    \item[$\mathrm{Weight}$] Weight (kg~head$^{-1}$)
    \item[$\mathrm{ADG}$] Average daily gain (kg~head$^{-1}$~day$^{-1}$), calculated
          using Eq.~3.2.1.6
    \item[$\mathrm{days}$] Total number of days since start of management period
    \item[$\mathrm{wt}_{\mathrm{initial}}$] Initial weight (kg~head$^{-1}$) (Table~16,
          by cattle group)
\end{description}

\bigskip
\noindent\textbf{Net energy requirements}

\begin{equation}
    \mathrm{NE}_{\mathrm{maintenance}} = C_f \times \mathrm{Weight}^{0.75}
    \tag{Eq.\ 3.2.1.2}
\end{equation}

\noindent IPCC (2019), Eq.~10.3

\noindent where
\begin{description}[font=\normalfont, leftmargin=9em, style=sameline, nosep]
    \item[$\mathrm{NE}_{\mathrm{maintenance}}$] Net energy for maintenance
          (MJ~head$^{-1}$~day$^{-1}$)
    \item[$C_f$] Maintenance coefficient (MJ~day$^{-1}$~kg$^{-1}$) (Table~16,
          by dairy cattle group)
\end{description}

\begin{equation}
    \mathrm{NE}_{\mathrm{activity}} = C_a \times \mathrm{NE}_{\mathrm{maintenance}}
    \tag{Eq.\ 3.2.1.3}
\end{equation}

\noindent IPCC (2019), Eq.~10.4

\noindent where
\begin{description}[font=\normalfont, leftmargin=9em, style=sameline, nosep]
    \item[$\mathrm{NE}_{\mathrm{activity}}$] Net energy for activity
          (MJ~head$^{-1}$~day$^{-1}$)
    \item[$C_a$] Feeding activity coefficient (Table~17, by activity type)
\end{description}

\noindent For lactating dairy cows only (use only when cows are lactating --- as a
default, Holos assumes 305 days lactation per year):

\begin{equation}
    \mathrm{NE}_{\mathrm{lactation}} =
        \mathrm{milk}_{\mathrm{prod}} \times
        \bigl(1.47 + 0.40 \times \mathrm{fat}_{\mathrm{content}}\bigr)
    \tag{Eq.\ 3.2.1.4}
\end{equation}

\noindent IPCC (2019), Eq.~10.8

\noindent where
\begin{description}[font=\normalfont, leftmargin=9em, style=sameline, nosep]
    \item[$\mathrm{NE}_{\mathrm{lactation}}$] Net energy for lactation
          (MJ~head$^{-1}$~day$^{-1}$)
    \item[$\mathrm{milk}_{\mathrm{prod}}$] Milk production (kg~head$^{-1}$~day$^{-1}$),
          by province (Table~21)
    \item[$\mathrm{fat}_{\mathrm{content}}$] Fat content (\%). A default value of 3.71\%
          is used (Canadian Dairy Information Centre)
\end{description}

\noindent\textit{Note: fat\_content is entered as a percentage (e.g.\ as 4 not 0.04).}

\noindent For pregnant dairy cows only:

\begin{equation}
    \mathrm{NE}_{\mathrm{pregnancy}} = 0.10 \times \mathrm{NE}_{\mathrm{maintenance}}
    \tag{Eq.\ 3.2.1.5}
\end{equation}

\noindent IPCC (2019), Eq.~10.13

\noindent where
\begin{description}[font=\normalfont, leftmargin=9em, style=sameline, nosep]
    \item[$\mathrm{NE}_{\mathrm{pregnancy}}$] Net energy for pregnancy
          (MJ~head$^{-1}$~day$^{-1}$)
    \item[$0.1$] Pregnancy coefficient --- IPCC (2019) default for cattle and buffalo
\end{description}

\noindent\textit{This equation averages pregnancy energy requirements over the
entire year.}

\bigskip
\noindent\textbf{Average daily gain, net energy for gain}

\begin{equation}
    \mathrm{ADG} = \frac{\mathrm{wt}_{\mathrm{final}} - \mathrm{wt}_{\mathrm{initial}}}
                        {\#\mathrm{days}}
    \tag{Eq.\ 3.2.1.6}
\end{equation}

\noindent where
\begin{description}[font=\normalfont, leftmargin=9em, style=sameline, nosep]
    \item[$\mathrm{ADG}$] Average daily gain (kg~head$^{-1}$~day$^{-1}$)
    \item[$\mathrm{wt}_{\mathrm{initial}}$] Initial weight (kg~head$^{-1}$) (Table~16,
          by dairy cattle group)
    \item[$\mathrm{wt}_{\mathrm{final}}$] Final weight (kg~head$^{-1}$) (Table~16,
          by dairy cattle group)
    \item[$\#\mathrm{days}$] Number of days --- includes all months where $>0$ animals
          are entered
\end{description}

\begin{equation}
    \mathrm{NE}_{\mathrm{gain}} = 22.02 \times
        \frac{\mathrm{Weight}}{C_d \times \mathrm{wt}_{\mathrm{final,milkcow}}^{0.75}}
        \times \mathrm{ADG}^{1.097}
    \tag{Eq.\ 3.2.1.7}
\end{equation}

\noindent IPCC (2019), Eq.~10.6

\noindent where
\begin{description}[font=\normalfont, leftmargin=9em, style=sameline, nosep]
    \item[$\mathrm{NE}_{\mathrm{gain}}$] Net energy for gain (MJ~head$^{-1}$~day$^{-1}$)
    \item[$\mathrm{Weight}$] Weight (kg~head$^{-1}$)
    \item[$\mathrm{wt}_{\mathrm{final,milkcow}}$] Final weight of animal (kg)
    \item[$C_d$] Gain coefficient (Table~16, by cattle group)
\end{description}

\bigskip
\noindent\textbf{Ratios of net energy available to digestible energy}

\begin{equation}
    \mathrm{REM} = 1.123 - 4.092 \times 10^{-3} \times \mathrm{TDN}
        + 1.126 \times 10^{-5} \times \mathrm{TDN}^{2}
        - \frac{25.4}{\mathrm{TDN}}
    \tag{Eq.\ 3.2.1.8}
\end{equation}

\noindent IPCC (2019), Eq.~10.14

\noindent where
\begin{description}[font=\normalfont, leftmargin=6em, style=sameline, nosep]
    \item[$\mathrm{REM}$] Ratio of net energy available in diet for maintenance to
          digestible energy consumed
    \item[$\mathrm{TDN}$] Percent total digestible nutrients in feed (Table~18,
          by diet)
\end{description}

\noindent\textit{Note: TDN is to be entered as a percentage (e.g.\ as 81 not 0.81).}

\begin{equation}
    \mathrm{REG} = 1.164 - 5.160 \times 10^{-3} \times \mathrm{TDN}
        + 1.308 \times 10^{-5} \times \mathrm{TDN}^{2}
        - \frac{37.44}{\mathrm{TDN}}
    \tag{Eq.\ 3.2.1.9}
\end{equation}

\noindent IPCC (2019), Eq.~10.15

\noindent where
\begin{description}[font=\normalfont, leftmargin=6em, style=sameline, nosep]
    \item[$\mathrm{REG}$] Ratio of net energy available in diet for gain to
          digestible energy consumed
\end{description}

\noindent\textit{Note: TDN is to be entered as a percentage (e.g.\ as 81 not 0.81).}

\bigskip
\noindent\textbf{Gross energy intake}

\begin{multline}
    \mathrm{GEI} =
        \left(\frac{\mathrm{NE}_{\mathrm{maintenance}} + \mathrm{NE}_{\mathrm{activity}}
            + \mathrm{NE}_{\mathrm{lactation}} + \mathrm{NE}_{\mathrm{pregnancy}}}
            {\mathrm{REM}}
        + \frac{\mathrm{NE}_{\mathrm{gain}}}{\mathrm{REG}}\right)
        \times \frac{100}{\mathrm{TDN}}
    \tag{Eq.\ 3.2.1.10}
\end{multline}

\noindent IPCC (2019), Eq.~10.16

\noindent where
\begin{description}[font=\normalfont, leftmargin=6em, style=sameline, nosep]
    \item[$\mathrm{GEI}$] Gross energy intake (MJ~head$^{-1}$~day$^{-1}$)
\end{description}

\noindent\textit{Note: TDN is to be entered as a percentage (e.g.\ as 81 not 0.81).}

\bigskip
\noindent\textbf{Enteric CH$_4$ emissions}

\begin{equation}
    \mathrm{CH}_{4,\mathrm{enteric\text{-}rate}} =
        \mathrm{GEI} \times \frac{Y_m}{55.65}
        \times \left(1 - \frac{\mathrm{AR}}{100}\right)
    \tag{Eq.\ 3.2.1.11}
\end{equation}

\noindent Derived from IPCC (2019), Eq.~10.21

\noindent where
\begin{description}[font=\normalfont, leftmargin=9em, style=sameline, nosep]
    \item[$\mathrm{CH}_{4,\mathrm{enteric\text{-}rate}}$] Enteric CH$_4$ emission rate
          (kg~head$^{-1}$~day$^{-1}$)
    \item[$Y_m$] CH$_4$ conversion factor (Table~18, by diet)
    \item[$55.65$] Energy content of CH$_4$ (MJ~kg$^{-1}$~CH$_4$)
    \item[$\mathrm{AR}$] Additive reduction factor (Table~19, by additive)
\end{description}

\begin{equation}
    \mathrm{CH}_{4,\mathrm{enteric}} =
        \mathrm{CH}_{4,\mathrm{enteric\text{-}rate}} \times \#\mathrm{cattle}
    \tag{Eq.\ 3.2.1.12}
\end{equation}

\noindent where
\begin{description}[font=\normalfont, leftmargin=9em, style=sameline, nosep]
    \item[$\mathrm{CH}_{4,\mathrm{enteric}}$] Enteric CH$_4$ emissions by animal group
          (kg~CH$_4$~day$^{-1}$)
    \item[$\#\mathrm{cattle}$] Number of dairy cattle
\end{description}

\noindent\textit{For dairy calves:} In Holos, a single default management period for
pre-weaned dairy calves is assumed, during which time these animals are exclusively
milk-fed and enteric CH$_4$ emissions are assumed to be zero. If the model user wishes
to include a weaned dairy calf group, they should do this by adding another management
period under the `Dairy heifers' animal group to model emissions from older calves eating
a partial or full forage diet. To estimate the daily Weight of milk-fed dairy calves,
Holos uses the same approach used for beef calves --- see Eq.~3.1.2.1. For dairy calves
on a milk-only diet, the DMI is estimated as: $\mathrm{DMI} = \mathrm{Weight} \times
0.01$. The GEI is also estimated using the beef calf approach --- see Eq.~3.1.2.3.

\bigskip
\noindent\textbf{Other enteric CH$_4$ estimates}

In addition to using the default IPCC (2019) $Y_m$ conversion factors to estimate
enteric CH$_4$ emissions from dairy cattle, the inclusion of different mathematical
equations (linear and exponential) with varying input requirements provides the model
user (and in particular researchers) with additional options. For dairy cattle, the
selected equations were developed and evaluated by Niu et al.\ (2018) and Benaouda
et al.\ (2019). Four equations were selected and divided into two groups based on the
concentration of forage in the diet. Two equations are for diets containing $\geq 25$\%
forage DM (high forage diets, Eq.~3.2.1.13 and Eq.~3.2.1.14) and the remaining two
equations are for diets containing $< 25$\% forage DM (low forage diets,
Eq.~3.2.1.15 and Eq.~3.2.1.16).

\noindent\textit{High forage diets alternatives ($\geq 25$\% dietary forage)}

\begin{equation}
    \mathrm{CH}_{4,\mathrm{enteric}(\mathrm{D1})\mathrm{\text{-}rate}} =
        \left(20 + 35.8 \times \mathrm{DMI}
            - 0.5 \times \mathrm{DMI}^{2}\right)
        \times 0.714 \times \frac{1}{1000}
    \tag{Eq.\ 3.2.1.13}
\end{equation}

\noindent Ramin and Huhtanen (2013)

\noindent where
\begin{description}[font=\normalfont, leftmargin=9em, style=sameline, nosep]
    \item[$\mathrm{CH}_{4,\mathrm{enteric}(\mathrm{D1})\mathrm{\text{-}rate}}$]
          Enteric CH$_4$ emissions rate (kg~head$^{-1}$~day$^{-1}$)
    \item[$\mathrm{DMI}$] DM intake (kg~head$^{-1}$~day$^{-1}$), calculated using
          Eq.~11.3.1.1
\end{description}

\begin{equation}
    \mathrm{CH}_{4,\mathrm{enteric}(\mathrm{D2})\mathrm{\text{-}rate}} =
        \left[56.27 - 56.27 \times e^{-0.028 \times \mathrm{DMI}}\right]
        \times \frac{1}{55.65}
    \tag{Eq.\ 3.2.1.14}
\end{equation}

\noindent Mills et al.\ (2003), Eq.~NL1

\noindent where
\begin{description}[font=\normalfont, leftmargin=9em, style=sameline, nosep]
    \item[$\mathrm{CH}_{4,\mathrm{enteric}(\mathrm{D2})\mathrm{\text{-}rate}}$]
          Enteric CH$_4$ emissions rate (kg~head$^{-1}$~day$^{-1}$)
    \item[$\mathrm{DMI}$] DM intake (kg~head$^{-1}$~day$^{-1}$), calculated using
          Eq.~11.3.1.1
\end{description}

\noindent\textit{Low forage diets alternatives ($< 25$\% dietary forage)}

\begin{equation}
    \mathrm{CH}_{4,\mathrm{enteric}(\mathrm{D3})\mathrm{\text{-}rate}} =
        \left(2.16 + 0.493 \times \mathrm{DMI}
            - 1.36 \times \mathrm{ADFI}
            + 1.97 \times \mathrm{NDFI}\right)
        \times \frac{1}{55.65}
    \tag{Eq.\ 3.2.1.15}
\end{equation}

\noindent Ellis et al.\ (2007), Eq.~8d

\noindent where
\begin{description}[font=\normalfont, leftmargin=9em, style=sameline, nosep]
    \item[$\mathrm{CH}_{4,\mathrm{enteric}(\mathrm{D3})\mathrm{\text{-}rate}}$]
          Enteric CH$_4$ emissions rate (kg~head$^{-1}$~day$^{-1}$)
    \item[$\mathrm{DMI}$] DM intake (kg~head$^{-1}$~day$^{-1}$), calculated using
          Eq.~11.3.1.1
    \item[$\mathrm{ADFI}$] Acid detergent fiber intake (kg~head$^{-1}$~day$^{-1}$)
    \item[$\mathrm{NDFI}$] Neutral detergent fiber intake (kg~head$^{-1}$~day$^{-1}$)
\end{description}

\begin{equation}
    \mathrm{CH}_{4,\mathrm{enteric}(\mathrm{D4})\mathrm{\text{-}rate}} =
        \left(76.0 + 13.5 \times \mathrm{DMI}
            - 9.55 \times \mathrm{EE}
            + 2.24 \times \mathrm{NDF}\right)
        \times \frac{1}{1000}
    \tag{Eq.\ 3.2.1.16}
\end{equation}

\noindent Niu et al.\ (2018), Eq.~5

\noindent where
\begin{description}[font=\normalfont, leftmargin=9em, style=sameline, nosep]
    \item[$\mathrm{CH}_{4,\mathrm{enteric}(\mathrm{D4})\mathrm{\text{-}rate}}$]
          Enteric CH$_4$ emissions rate (kg~head$^{-1}$~day$^{-1}$)
    \item[$\mathrm{DMI}$] DM intake (kg~head$^{-1}$~day$^{-1}$), calculated using
          Eq.~11.3.1.1
    \item[$\mathrm{EE}$] Dietary fat/ether extract (\% of DM)
    \item[$\mathrm{NDF}$] Dietary neutral detergent fiber (\% of DM)
\end{description}

\noindent Holos V4 uses the following calculations:
$\mathrm{ADFI} = \mathrm{DMI}\,(\text{kg~d}^{-1}) \times [\mathrm{ADF\ in\ the\ diet\ or\ feed}\ (\%\mathrm{DM})/100]$;
$\mathrm{NDFI} = \mathrm{DMI}\,(\text{kg~d}^{-1}) \times [\mathrm{NDF\ in\ the\ diet\ or\ feed}\ (\%\mathrm{DM})/100]$

\begin{equation}
    \mathrm{CH}_{4,\mathrm{enteric}(n)} =
        \mathrm{CH}_{4,\mathrm{enteric}(n)\mathrm{\text{-}rate}} \times \#\mathrm{cattle}
    \tag{Eq.\ 3.2.1.17}
\end{equation}

\noindent where
\begin{description}[font=\normalfont, leftmargin=9em, style=sameline, nosep]
    \item[$\mathrm{CH}_{4,\mathrm{enteric}(n)}$] Enteric CH$_4$ emissions
          (kg~CH$_4$), $n = \mathrm{D1, D2, D3, D4}$
    \item[$\#\mathrm{cattle}$] Number of cattle
\end{description}

% ·····················································
\subsection{Enteric CH\texorpdfstring{$_4$}{4} emissions from dairy calves}
\label{subsec:dairy-calves}

For dairy calves, Eq.~3.2.2.1 is used to calculate enteric CH$_4$ emissions from
milk-fed dairy calves.

\begin{equation}
    \mathrm{CH}_{4,\mathrm{enteric}} = 0
    \tag{Eq.\ 3.2.2.1}
\end{equation}

\noindent IPCC (2019)

\noindent where
\begin{description}[font=\normalfont, leftmargin=9em, style=sameline, nosep]
    \item[$\mathrm{CH}_{4,\mathrm{enteric}}$] Enteric CH$_4$ emissions from dairy
          calves (kg~CH$_4$~day$^{-1}$)
\end{description}

\noindent IPCC (2019) states that ``A CH$_4$ conversion factor of zero is assumed for
all juveniles consuming only milk (i.e., milk-fed lambs as well as calves).'' Therefore,
there are no enteric CH$_4$ emissions associated with milk-fed dairy calves.

% ·····················································
\subsection{Total enteric CH\texorpdfstring{$_4$}{4} emissions}
\label{subsec:dairy-total}

Emissions should be summed for all relevant management periods and dairy cattle groups.

\begin{equation}
    \mathrm{Total\_CH}_{4,\mathrm{enteric}} =
        \sum_{\mathrm{all\ scenario\ cattle}}
        \mathrm{CH}_{4,\mathrm{enteric}}
    \tag{Eq.\ 3.2.3.1}
\end{equation}

\noindent IPCC (2019), Eq.~10.20

\noindent where
\begin{description}[font=\normalfont, leftmargin=9em, style=sameline, nosep]
    \item[$\mathrm{Total\_CH}_{4,\mathrm{enteric}}$] Total enteric CH$_4$ emissions
          from dairy cattle (kg~CH$_4$~year$^{-1}$)
    \item[$\mathrm{CH}_{4,\mathrm{enteric}}$] Enteric CH$_4$ emissions (kg~CH$_4$),
          by dairy cattle group
\end{description}

% ─────────────────────────────────────────────────────────────────────────────
\section{Sheep}
\label{sec:sheep}

Management data are input for each management period. Emissions are calculated
on a daily basis.

\noindent\textbf{Assumptions:}
\begin{itemize}[leftmargin=1.5em]
    \item The user specifies the number of ewes and number of lambs, from which the
          number of lambs per ewe (lamb:ewe ratio) is calculated
    \item The pregnancy twinning rate is based on the current lamb:ewe ratio
          (assumed static)
    \item There are no emissions from nursing lambs
    \item Sheep feed intake is equal to energy requirements
    \item The sex ratio of the feedlot is 1:1
    \item All feed is utilized --- waste feed emissions are not accounted for
    \item All barn manure is handled as solid storage (stockpiled) manure and
          land-applied at least once per year
\end{itemize}

% ·····················································
\subsection{Enteric CH\texorpdfstring{$_4$}{4} emissions from sheep}
\label{subsec:sheep-enteric}

Enteric CH$_4$ calculations should be completed for each sheep group.

\begin{equation}
    \mathrm{Weight} = \mathrm{ADG} \times \mathrm{days} + \mathrm{wt}_{\mathrm{initial}}
    \tag{Eq.\ 3.3.1.1}
\end{equation}

\noindent where
\begin{description}[font=\normalfont, leftmargin=9em, style=sameline, nosep]
    \item[$\mathrm{Weight}$] Weight (kg~head$^{-1}$)
    \item[$\mathrm{ADG}$] Average daily gain (kg~head$^{-1}$~day$^{-1}$), calculated
          using Eq.~3.3.1.9
    \item[$\mathrm{days}$] Total number of days since start of management period
    \item[$\mathrm{wt}_{\mathrm{initial}}$] Initial weight (kg~head$^{-1}$) (Table~22,
          by sheep group)
\end{description}

\bigskip
\noindent\textbf{Lamb:ewe ratio}

This is calculated only for those days when there are lambs with ewes; the average
yearly lamb:ewe ratio is used.

\begin{equation}
    \mathrm{Lamb\_Ratio} = \frac{\#\mathrm{lambs}}{\#\mathrm{ewes}}
    \tag{Eq.\ 3.3.1.2}
\end{equation}

\noindent where
\begin{description}[font=\normalfont, leftmargin=9em, style=sameline, nosep]
    \item[$\mathrm{Lamb\_Ratio}$] Lamb:ewe ratio
    \item[$\#\mathrm{lambs}$] Number of lambs
    \item[$\#\mathrm{ewes}$] Number of ewes
\end{description}

\bigskip
\noindent\textbf{Net energy requirements}

\begin{equation}
    \mathrm{NE}_{\mathrm{maintenance}} = C_f \times \mathrm{Weight}^{0.75}
    \tag{Eq.\ 3.3.1.3}
\end{equation}

\noindent IPCC (2019), Eq.~10.3

\noindent where
\begin{description}[font=\normalfont, leftmargin=9em, style=sameline, nosep]
    \item[$\mathrm{NE}_{\mathrm{maintenance}}$] Net energy for maintenance
          (MJ~head$^{-1}$~day$^{-1}$)
    \item[$C_f$] Maintenance coefficient (MJ~day$^{-1}$~kg$^{-1}$) (Table~22,
          by sheep group)
\end{description}

\begin{equation}
    \mathrm{NE}_{\mathrm{activity}} = C_a \times \mathrm{Weight}
    \tag{Eq.\ 3.3.1.4}
\end{equation}

\noindent IPCC (2019), Eq.~10.5

\noindent where
\begin{description}[font=\normalfont, leftmargin=9em, style=sameline, nosep]
    \item[$\mathrm{NE}_{\mathrm{activity}}$] Net energy for activity
          (MJ~head$^{-1}$~day$^{-1}$)
    \item[$C_a$] Feeding activity coefficient (MJ~kg$^{-1}$) (Table~23,
          by activity type)
\end{description}

\noindent For lactating ewes only (use only when ewes are lactating and milk
production is \textbf{unknown} and when \#lambs $>$ 0):

\begin{equation}
    \mathrm{NE}_{\mathrm{lactation}} =
        5 \times \mathrm{WG}_{\mathrm{lamb}} \times \mathrm{EV}_{\mathrm{milk}}
    \tag{Eq.\ 3.3.1.5}
\end{equation}

\noindent Derived from IPCC (2019), Eq.~10.10

\noindent where
\begin{description}[font=\normalfont, leftmargin=9em, style=sameline, nosep]
    \item[$\mathrm{NE}_{\mathrm{lactation}}$] Net energy for lactation
          (MJ~head$^{-1}$~day$^{-1}$)
    \item[$\mathrm{WG}_{\mathrm{lamb}}$] Daily pre-weaning weight gain of lamb(s)
          (Table~24, by Lamb:ewe ratio)
    \item[$5$] When milk production is not known, AFRC (1990) indicates that for a
          single birth, the milk yield is about 5 times the weight gain of the lamb.
          For multiple births, the total annual milk production can be estimated as
          five times the increase in combined weight gain of all lambs birthed by a
          single ewe.
    \item[$\mathrm{EV}_{\mathrm{milk}}$] Energy required to produce 1\,kg of milk
          (MJ~kg$^{-1}$)
\end{description}

\noindent Holos V4 uses the following constant value:
$\mathrm{EV}_{\mathrm{milk}} = 4.6$\,MJ~kg$^{-1}$ (IPCC 2019). The default IPCC (2019)
$\mathrm{EV}_{\mathrm{milk}}$ value of 4.6\,MJ~kg$^{-1}$ can be used for sheep
(AFRC 1993; AFRC 1995), which corresponds to a milk fat content of 7\% by weight.

\noindent For lactating ewes only (use only when ewes are lactating and milk
production is \textbf{known} and when \#lambs $>$ 0):

\begin{equation}
    \mathrm{NE}_{\mathrm{lactation}} =
        \mathrm{Milk} \times \mathrm{EV}_{\mathrm{milk}}
    \tag{Eq.\ 3.3.1.6}
\end{equation}

\noindent IPCC (2019), Eq.~10.9

\noindent where
\begin{description}[font=\normalfont, leftmargin=9em, style=sameline, nosep]
    \item[$\mathrm{Milk}$] Amount of milk produced (kg~head$^{-1}$~day$^{-1}$).
          A default value of 2\,kg~head$^{-1}$~day$^{-1}$ is used (based on a value
          of approximately 2\,L per sheep per day (Farm and Food Care Ontario 2016)),
          assuming that 1\,L of milk weighs approx.\ 1\,kg.
\end{description}

\noindent For pregnant ewes only:

\begin{equation}
    \mathrm{NE}_{\mathrm{pregnancy}} =
        C_{\mathrm{preg}} \times \mathrm{NE}_{\mathrm{maintenance}}
    \tag{Eq.\ 3.3.1.7}
\end{equation}

\noindent IPCC (2019), Eq.~10.13

\noindent where
\begin{description}[font=\normalfont, leftmargin=9em, style=sameline, nosep]
    \item[$\mathrm{NE}_{\mathrm{pregnancy}}$] Net energy for pregnancy
          (MJ~head$^{-1}$~day$^{-1}$)
    \item[$C_{\mathrm{preg}}$] Pregnancy coefficient (Table~25, by Lamb:ewe ratio)
\end{description}

\noindent\textit{This equation averages pregnancy energy requirements over the
entire year.}

\noindent For ewes and rams only:

\begin{equation}
    \mathrm{NE}_{\mathrm{wool}} =
        \frac{\mathrm{EV}_{\mathrm{wool}} \times \mathrm{wool}_{\mathrm{prod}}}{365}
    \tag{Eq.\ 3.3.1.8}
\end{equation}

\noindent IPCC (2019), Eq.~10.12

\noindent where
\begin{description}[font=\normalfont, leftmargin=9em, style=sameline, nosep]
    \item[$\mathrm{NE}_{\mathrm{wool}}$] Net energy for wool production
          (MJ~head$^{-1}$~day$^{-1}$)
    \item[$\mathrm{EV}_{\mathrm{wool}}$] Energy value of 1\,kg of wool (MJ~kg$^{-1}$)
    \item[$\mathrm{wool}_{\mathrm{prod}}$] Wool production (kg~year$^{-1}$) (Table~22,
          by sheep group)
    \item[$365$] Number of days in year
\end{description}

\noindent Holos V4 uses the following constant value:
$\mathrm{EV}_{\mathrm{wool}} = 24$\,MJ~kg$^{-1}$ (IPCC 2019).

\bigskip
\noindent\textbf{Average daily gain, net energy for gain}

ADG is not entered into the model interface for ewes and rams, but is required for
sheep feedlot animals. For ewes and rams:

\begin{equation}
    \mathrm{ADG} = \frac{\mathrm{wt}_{\mathrm{final}} - \mathrm{wt}_{\mathrm{initial}}}
                        {\#\mathrm{days}}
    \tag{Eq.\ 3.3.1.9}
\end{equation}

\noindent where
\begin{description}[font=\normalfont, leftmargin=9em, style=sameline, nosep]
    \item[$\mathrm{ADG}$] Average daily gain (kg~head$^{-1}$~day$^{-1}$)
    \item[$\mathrm{wt}_{\mathrm{initial}}$] Initial weight (kg~head$^{-1}$) (Table~22,
          by sheep group)
    \item[$\mathrm{wt}_{\mathrm{final}}$] Final weight (kg~head$^{-1}$) (Table~22,
          by sheep group)
    \item[$\#\mathrm{days}$] Number of days in management period --- includes all days
          where animals $> 0$ are entered
\end{description}

\noindent For feedlot sheep (weaned lambs), ewes and rams:

\begin{equation}
    \mathrm{NE}_{\mathrm{gain}} =
        \mathrm{ADG} \times \bigl(a + b \times \mathrm{Weight}\bigr)
    \tag{Eq.\ 3.3.1.10}
\end{equation}

\noindent Based on IPCC (2019), Eq.~10.7

\noindent where
\begin{description}[font=\normalfont, leftmargin=9em, style=sameline, nosep]
    \item[$\mathrm{NE}_{\mathrm{gain}}$] Net energy for gain (MJ~head$^{-1}$~day$^{-1}$)
    \item[$a$] Coefficient $a$ (MJ~kg$^{-1}$) (Table~22, by sheep group)
    \item[$b$] Coefficient $b$ (MJ~kg$^{-2}$) (Table~22, by sheep group)
    \item[$\mathrm{Weight}$] Weight (kg~head$^{-1}$)
    \item[$\mathrm{ADG}$] Average daily gain (kg~head$^{-1}$~day$^{-1}$), calculated
          using Eq.~3.3.1.9
\end{description}

\bigskip
\noindent\textbf{Ratios of net energy available to digestible energy}

\begin{equation}
    \mathrm{REM} = 1.123 - 4.092 \times 10^{-3} \times \mathrm{TDN}
        + 1.126 \times 10^{-5} \times \mathrm{TDN}^{2}
        - \frac{25.4}{\mathrm{TDN}}
    \tag{Eq.\ 3.3.1.11}
\end{equation}

\noindent IPCC (2019), Eq.~10.14

\noindent where
\begin{description}[font=\normalfont, leftmargin=6em, style=sameline, nosep]
    \item[$\mathrm{REM}$] Ratio of net energy available in diet for maintenance to
          digestible energy consumed
    \item[$\mathrm{TDN}$] Percent digestible energy in feed (Table~26, by diet)
\end{description}

\noindent\textit{Note: TDN is to be entered as a percentage (e.g.\ as 81 not 0.81).}

\begin{equation}
    \mathrm{REG} = 1.164 - 5.16 \times 10^{-3} \times \mathrm{TDN}
        + 1.308 \times 10^{-5} \times \mathrm{TDN}^{2}
        - \frac{37.4}{\mathrm{TDN}}
    \tag{Eq.\ 3.3.1.12}
\end{equation}

\noindent IPCC (2019), Eq.~10.15

\noindent where
\begin{description}[font=\normalfont, leftmargin=6em, style=sameline, nosep]
    \item[$\mathrm{REG}$] Ratio of net energy available in diet for gain to
          digestible energy consumed
\end{description}

\noindent\textit{Note: TDN is to be entered as a percentage (e.g.\ as 81 not 0.81).}

\bigskip
\noindent\textbf{Gross energy intake}

\begin{multline}
    \mathrm{GEI} =
        \left(\frac{\mathrm{NE}_{\mathrm{maintenance}} + \mathrm{NE}_{\mathrm{activity}}
            + \mathrm{NE}_{\mathrm{lactation}} + \mathrm{NE}_{\mathrm{pregnancy}}}
            {\mathrm{REM}}
        + \frac{\mathrm{NE}_{\mathrm{gain}} + \mathrm{NE}_{\mathrm{wool}}}
               {\mathrm{REG}}\right)
        \times \frac{100}{\mathrm{TDN}}
    \tag{Eq.\ 3.3.1.13}
\end{multline}

\noindent IPCC (2019), Eq.~10.16

\noindent where
\begin{description}[font=\normalfont, leftmargin=6em, style=sameline, nosep]
    \item[$\mathrm{GEI}$] Gross energy intake (MJ~head$^{-1}$~day$^{-1}$)
\end{description}

\noindent\textit{Note: TDN is to be entered as a percentage (e.g.\ as 81 not 0.81).}

\bigskip
\noindent\textbf{Enteric CH$_4$ emissions}

\begin{equation}
    \mathrm{CH}_{4,\mathrm{enteric\text{-}rate}} =
        \mathrm{GEI} \times \frac{Y_m}{55.65}
    \tag{Eq.\ 3.3.1.14}
\end{equation}

\noindent IPCC (2019), Eq.~10.21

\noindent where
\begin{description}[font=\normalfont, leftmargin=9em, style=sameline, nosep]
    \item[$\mathrm{CH}_{4,\mathrm{enteric\text{-}rate}}$] Enteric CH$_4$ emission
          rate (kg~head$^{-1}$~day$^{-1}$)
    \item[$Y_m$] CH$_4$ conversion factor (Table~26, by sheep group)
    \item[$55.65$] Energy content of CH$_4$ (MJ~kg$^{-1}$~CH$_4$)
\end{description}

\begin{equation}
    \mathrm{CH}_{4,\mathrm{enteric}} =
        \mathrm{CH}_{4,\mathrm{enteric\text{-}rate}} \times \#\mathrm{sheep}
    \tag{Eq.\ 3.3.1.15}
\end{equation}

\noindent where
\begin{description}[font=\normalfont, leftmargin=9em, style=sameline, nosep]
    \item[$\mathrm{CH}_{4,\mathrm{enteric}}$] Enteric CH$_4$ emissions, by sheep
          group (kg~CH$_4$~day$^{-1}$)
    \item[$\#\mathrm{sheep}$] Number of sheep
\end{description}

% ·····················································
\subsection{Total enteric CH\texorpdfstring{$_4$}{4} emissions}
\label{subsec:sheep-total}

Emissions should be summed for all relevant management periods and sheep groups.

\begin{equation}
    \mathrm{Total\_CH}_{4,\mathrm{enteric}} =
        \sum_{\mathrm{all\ scenario\ sheep}}
        \mathrm{CH}_{4,\mathrm{enteric}}
    \tag{Eq.\ 3.3.2.1}
\end{equation}

\noindent IPCC (2019), Eq.~10.20

\noindent where
\begin{description}[font=\normalfont, leftmargin=9em, style=sameline, nosep]
    \item[$\mathrm{Total\_CH}_{4,\mathrm{enteric}}$] Total enteric CH$_4$ emissions
          from sheep (kg~CH$_4$~year$^{-1}$)
    \item[$\mathrm{CH}_{4,\mathrm{enteric}}$] Enteric CH$_4$ emissions (kg~CH$_4$)
\end{description}

% ─────────────────────────────────────────────────────────────────────────────
\section{Swine, Poultry and other livestock}
\label{sec:swine-poultry}

Management data are input for each management period. Emissions are calculated
on a daily basis.

\noindent\textbf{Assumptions:}
\begin{itemize}[leftmargin=1.5em]
    \item All feed is utilized --- waste feed emissions are not accounted for
    \item All barn manure uses the same handling system
    \item Manure must be applied at least once per year
    \item There are no emissions from nursing piglets or other unweaned young
    \item There are no emissions from chicks or poults
\end{itemize}

% ·····················································
\subsection{Enteric CH\texorpdfstring{$_4$}{4} emissions from swine, poultry and
other livestock}
\label{subsec:swine-enteric}

Daily enteric CH$_4$ calculations should be completed for each animal group within
each livestock type.

\begin{equation}
    \mathrm{CH}_{4,\mathrm{enteric}} =
        \frac{\mathrm{CH}_{4,\mathrm{enteric\text{-}rate}}}{365}
        \times \#\mathrm{animals}
    \tag{Eq.\ 3.4.1.1}
\end{equation}

\noindent IPCC (2019), Eq.~10.19

\noindent where
\begin{description}[font=\normalfont, leftmargin=9em, style=sameline, nosep]
    \item[$\mathrm{CH}_{4,\mathrm{enteric}}$] Enteric CH$_4$ emissions
          (kg~CH$_4$~day$^{-1}$)
    \item[$\mathrm{CH}_{4,\mathrm{enteric\text{-}rate}}$] Yearly enteric CH$_4$
          emission rate (Verg\'{e} et al.\ (2009) for swine; IPCC (2019) for all
          other livestock) (Table~27)
    \item[$365$] Number of days in year
\end{description}

% ·····················································
\subsection{Total enteric CH\texorpdfstring{$_4$}{4} emissions}
\label{subsec:swine-total}

Emissions should be summed for all management periods and animal groups within each
livestock type.

\begin{equation}
    \mathrm{Total\_CH}_{4,\mathrm{enteric}} =
        \sum_{\mathrm{all\ scenario\ animals}}
        \mathrm{CH}_{4,\mathrm{enteric}}
    \tag{Eq.\ 3.4.2.1}
\end{equation}

\noindent IPCC (2019), Eq.~10.20

\noindent where
\begin{description}[font=\normalfont, leftmargin=9em, style=sameline, nosep]
    \item[$\mathrm{Total\_CH}_{4,\mathrm{enteric}}$] Total enteric CH$_4$ emissions
          (kg~CH$_4$~production cycle$^{-1}$) --- in the event that the production
          cycle equates to 365 days, this total is a yearly total
    \item[$\mathrm{CH}_{4,\mathrm{enteric}}$] Enteric CH$_4$ emissions (kg~CH$_4$),
          by animal group
\end{description}

% ─────────────────────────────────────────────────────────────────────────────
\section{Scaling up to yearly estimates}
\label{sec:scaling-up}

\noindent\textit{Note: Scaling up of the detailed emission report to annual values has
been disabled in V4, and all outputs reported in Holos are currently for a single
production cycle only. All Holos estimates, including those on the Manure Management
report, are for a single production cycle only --- if the model user wants to know these
estimates for an entire year (365 days), they will need to scale-up these estimates
themselves outside of the Holos interface --- please see below for the method to do so.}

\noindent\textbf{Scaling up of emissions estimates from production cycle to annual values}

In Canada, some livestock production systems are typically ``all in, all out'' systems,
where all animals enter and leave the barn at the same time with a number of days between
production periods to clean the facilities before the next group of animals arrives.
These systems include feedlot cattle (steers and heifers), swine, and poultry production.
Therefore, to calculate total annual emissions from a livestock operation where there are
multiple production cycles per year, the number of production days per year (excluding
non-productive days between production cycles) are considered. Daily emissions estimates
are multiplied by the total number of production days in a year.

In Holos, for each animal group, a single default production cycle (composed of one or
more management periods) is provided in the interface, detailing the number of production
days in each cycle. A default number of rest days between production cycles is also
provided (Table~28). For all animal groups not included in Table~28, we assume that
either their production cycle lasts 365 days (e.g., beef cows and bulls) or that they
typically have just a single production cycle per year (e.g., beef calves). This
`scaling-up' should be carried out in Holos for the following GHG emissions estimates
reported in the Detailed Emission Report on the Results screen: `Enteric CH$_4$',
`Manure CH$_4$', `Direct N$_2$O', `Indirect N$_2$O', `Farm energy CO$_2$',
`Upstream CO$_2$' and `Subtotal'.

\begin{equation}
    \mathrm{No.Prod.Cycles} =
        \frac{365}{\mathrm{No.Days}_{\mathrm{Cycle}}
              + \mathrm{No.Days}_{\mathrm{Rest}}}
    \tag{Eq.\ 3.5.1.1}
\end{equation}

\noindent where
\begin{description}[font=\normalfont, leftmargin=9em, style=sameline, nosep]
    \item[$\mathrm{No.Prod.Cycles}$] Total number of production cycles per year,
          by animal group
    \item[$\mathrm{No.Days}_{\mathrm{Cycle}}$] Number of days in a single production
          cycle, by animal group. This is calculated in Holos as the sum of days in
          all specified management periods for a given animal group; typical values for
          the average number of days in a single production cycle for certain animal
          groups are provided, for reference, in Table~28
    \item[$\mathrm{No.Days}_{\mathrm{Rest}}$] Number of non-production or rest days
          between production cycles, by animal group, user-specified or default
          (Table~28)
\end{description}

\begin{equation}
    \mathrm{CH}_{4,\mathrm{enteric,annual}} =
        \mathrm{CH}_{4,\mathrm{enteric}} \times \mathrm{No.Prod.Cycles}
    \tag{Eq.\ 3.5.1.2}
\end{equation}

\noindent where
\begin{description}[font=\normalfont, leftmargin=9em, style=sameline, nosep]
    \item[$\mathrm{CH}_{4,\mathrm{enteric,annual}}$] Total enteric CH$_4$ emissions
          for all production days in a year (kg~CH$_4$~year$^{-1}$), by animal group
    \item[$\mathrm{CH}_{4,\mathrm{enteric}}$] Enteric CH$_4$ emissions for all animals
          and all management periods within a single production cycle (kg~CH$_4$), by
          animal group
\end{description}

% ─────────────────────────────────────────────────────────────────────────────
\section{References}
\label{sec:enteric-refs}

\subsection*{3.1 Beef}
\begin{description}[leftmargin=0pt, style=nextline]

    \item[Ellis et al., 2009]
        Ellis, J.L., Kebreab, E., Odongo, N.E., Beauchemin, K., McGinn, S., Nkrumah,
        J.D., Moore, S.S., Christopherson, R., Murdoch, G.K., McBride, B.W., Okine,
        E.K., France, J., 2009. Modeling methane production from beef cattle using
        linear and nonlinear approaches. \textit{J.\ Anim.\ Sci.} 87, 1334--1345.
        \url{https://doi.org/10.2527/jas.2007-0725}

    \item[Escobar-Bahamondes et al., 2017]
        Escobar-Bahamondes, P., Oba, M., Beauchemin, K., 2017. Universally applicable
        methane prediction equations for beef cattle fed high- or low-forage diets.
        \textit{Can.\ J.\ Anim.\ Sci.} 94.
        \url{https://doi.org/10.1139/CJAS-2016-0042}

    \item[IPCC, 2019]
        2019 Refinement to the 2006 IPCC Guidelines for National Greenhouse Gas
        Inventories, Calvo Buendia, E., Tanabe, K., Kranjc, A., Baasansuren, J.,
        Fukuda, M., Ngarize S., Osako, A., Pyrozhenko, Y., Shermanau, P.\ and
        Federici, S.\ (eds). Published: IPCC, Switzerland.\\
        \url{https://www.ipcc.ch/report/2019-refinement-to-the-2006-ipcc-guidelines-for-national-greenhouse-gas-inventories/}

    \item[Legesse et al., 2016]
        Legesse G., Beauchemin K.A., Ominski K.H., McGeough E.J., Kroebel R.,
        MacDonald D., Little S.M., and McAllister T.A., 2016. Greenhouse gas emissions
        of Canadian beef production in 1981 as compared with 2011.
        \textit{Animal Production Science} 56, 153--168.
        \url{https://doi.org/10.1071/AN15386}

    \item[Legesse et al., 2018]
        Legesse, G., Cordeiro, M.R.C., Ominski, K.H., Beauchemin, K.A., Kroebel, R.,
        McGeough, E.J., Pogue, S., and McAllister, T.A., 2018. Water use intensity of
        Canadian beef production in 1981 as compared to 2011.
        \textit{Sci.\ Total Environ.} 619--620: 1030--1039.
        \url{https://doi.org/10.1016/j.scitotenv.2017.11.194}

    \item[van Lingen et al., 2019]
        van Lingen, H.J., Niu, M., Kebreab, E., et al., 2019. Prediction of enteric
        methane production, yield and intensity of beef cattle using an intercontinental
        database. \textit{Agriculture, Ecosystems and Environment} 283, 106575.
        \url{https://doi.org/10.1016/j.agee.2019.106575}

\end{description}

\subsection*{3.2 Dairy}
\begin{description}[leftmargin=0pt, style=nextline]

    \item[Benaouda et al., 2019]
        Benaouda, M., Martin, C., Li, X., Kebreab, E., et al., 2019. Evaluation of the
        performance of existing mathematical models predicting enteric methane emissions
        from ruminants: Animal categories and dietary mitigation strategies.
        \textit{Animal Feed Science and Technology} 255, 114207.
        \url{https://doi.org/10.1016/j.anifeedsci.2019.114207}

    \item[Ellis et al., 2007]
        Ellis, J.L., Kebreab, E., Odongo, N.E., McBride, B.W., Okine, E.K., France,
        J., 2007. Prediction of methane production from dairy and beef cattle.
        \textit{J.\ Dairy Sci.} 90, 3456--3466.
        \url{https://doi.org/10.3168/jds.2006-675}

    \item[IPCC, 2019]
        As cited in Section~3.1 Beef references.

    \item[Mills et al., 2003]
        Mills, J.A.N., Kebreab, E., Yates, C.M., Crompton, L.A., Cammell, S.B.,
        Dhanoa, M.S., Agnew, R.E., France, J., 2003. Alternative approaches to
        predicting methane emissions from dairy cows. \textit{J.\ Anim.\ Sci.} 81,
        3141--3150.
        \url{https://doi.org/10.2527/2003.81123141x}

    \item[Niu et al., 2018]
        Niu, M., Kebreab, E., Hristov, A.N., Oh, J., Arndt, C., et al., 2018.
        Prediction of enteric methane production, yield, and intensity in dairy cattle
        using an intercontinental database. \textit{Glob Change Biol.} 24: 3368--3389.
        \url{https://doi.org/10.1111/gcb.14094}

    \item[Ramin \& Huhtanen, 2013]
        Ramin, M. and Huhtanen, P., 2013. Development of equations for predicting
        methane emissions from ruminants. \textit{J Dairy Sci.} 96(4):2476--2493.
        \url{https://doi.org/10.3168/jds.2012-6095}

\end{description}

\subsection*{3.3 Sheep}
\begin{description}[leftmargin=0pt, style=nextline]

    \item[AFRC, 1990]
        Nutritive Requirements of Ruminant Animals: Energy. Rep.\ 5. Wallingford, UK:
        CAB International.

    \item[AFRC, 1993]
        Energy and Protein Requirements of Ruminants. An advisory manual prepared by
        the AFRC Technical Committee on Response to Nutrients. Wallingford, UK: CAB
        International, pp.\ 24--159.

    \item[AFRC, 1995]
        Energy and protein requirements of ruminants. An advisory manual prepared by
        the AFRC Technical Committee on Response to Nutrients. Wallingford, UK: CAB
        International, pp.\ 159.

    \item[Farm and Food Care Ontario, 2016]
        Facts and Figures about Canadian sheep.
        \url{https://www.farmfoodcareon.org/wp-content/uploads/2017/05/Fact-Sheet-Sheep-2016.pdf}

    \item[IPCC, 2019]
        As cited in Section~3.1 Beef references.

\end{description}

\subsection*{3.4 Swine, Poultry and others}
\begin{description}[leftmargin=0pt, style=nextline]

    \item[IPCC, 2019]
        As cited in Section~3.1 Beef references.

    \item[Verg\'{e} et al., 2009]
        Verg\'{e}, X.P.X., Dyer, J.A., Desjardins, R.L., Worth, D., 2009. Greenhouse
        gas emissions from the Canadian pork industry.
        \textit{Livestock Science} 121: 92--101.
        \url{https://doi.org/10.1016/j.livsci.2008.05.022}

\end{description}

\subsection*{3.5 Scaling up of emissions estimates from production cycle to annual values}
\begin{description}[leftmargin=0pt, style=nextline]

    \item[New-Life Mills, 2016]
        Pullet and Layer Management Guide: A Complete Guide to Profitable Performance.

    \item[Newlands et al., 2011]
        Newlands, N.K., Davidson, A., Howard, A., and Hill, H., 2011. Validation and
        inter-comparison of three methodologies for interpolating daily precipitation
        and temperature across Canada. \textit{Environmetrics} 22(2), 205--223.
        \url{https://doi.org/10.1002/env.1044}

    \item[Sheppard et al., 2009]
        Sheppard, S.C., Bittman, S., Beaulieu, M., and Sheppard, M.I., 2009. Ecoregion
        and farm-size differences in feed and manure nitrogen management: 1.\ Survey
        methods and results for poultry.
        \textit{Canadian Journal of Animal Science} 89: 1--19.
        \url{https://doi.org/10.4141/CJAS08054}

\end{description}
